\documentclass{beamer}

\usepackage[utf8]{inputenc}

\input{../helper}
\graphicspath{{../img/}}

\title{Introdução à Computação \\ análise de algorítmos (cont.)} 
\author{Alexandre Rademaker}
\date{}

\begin{document}

\begin{frame}
 \maketitle
\end{frame}

\begin{frame}[fragile,allowframebreaks]{análise de algorítmos}

  Como determinar a complexidade assintótica de um algoritmo?

  Vamos começar com alguns códigos simples.

  \framebreak

  \begin{lstlisting}[style=normalc]
    int main() {
      int a[n] = {...}, s = 0;
      for(int i = 0; i < n; i++) 
        s += a[i];
    }
  \end{lstlisting}

  Complexidade $O(n)$ passos, são $n$ instruções de soma elementares.

  \framebreak

  \begin{lstlisting}[style=normalc]
    int main() {
      int a[n] = {...}, s = 0;
      for(int i = 0; i < n; i++) 
        for(int j = 0; j < n; j++) 
          s += a[i] * a[j];
    }
  \end{lstlisting}

  Complexidade $O(n^2)$, desconsiderando fatores multiplicativos
  constantes (adição e multiplicação).

  \framebreak

  \begin{lstlisting}[style=normalc]
    int main() {
      int a[n] = {...}, s = 0;
      for(int i = 0; i < n; i++) {
        for(int j = 0; j < n; j++) 
          s += a[i] * a[j];
        for(int j = 0; j < n; j++) 
          for(int k = 0; k < n; k++) 
            s += a[i] - a[k];
      }
    }
  \end{lstlisting}

  Neste caso, $O(n (n + n^2)) = O(n^2 + n^3) = O(n^3)$.
  

  \framebreak

  \begin{lstlisting}[style=normalc]
    int main() {
      int a[n] = {...}, b[m] = {...}, s = 0;
      for(int i = 0; i < n; i++) 
        for(int j = 0; j < m; j++) 
          s += a[i] * b[j];
    }
  \end{lstlisting}

  Complexidade $O(nm)$, desconsiderando fatores multiplicativos
  constantes (adição e multiplicação).

  \framebreak

  \begin{lstlisting}[style=normalc]
    int main() {
      int a[n] = {...}, s = 0;
      for(int i = 0; i < n; i++) {
        if(a[i] > a[i+1])
          break;
      }
    }
  \end{lstlisting}
  
  Melhor caso? Pior caso? Caso médio?

  \framebreak

  \begin{lstlisting}[style=normalc]
    int func(int a[], n) {
      int s = 0;
      for(int i = 0; i < n; i++) 
        for(int j = 0; j < n; j++) 
          s += a[i] * a[j];
      return s;
    }
    
    int main() {
      int a[n] = {...};
      int s = func(a, n);
    }
  \end{lstlisting}
  
  Qual a complexidade? Qual a complexidade de + para valores grandes?
  Qual a complexidade de \ilcode{atoi} da \ilcode{stdlib.h}.
  
  
\end{frame}



\end{document}

%%% Local Variables:
%%% mode: latex
%%% TeX-master: t
%%% End:
